\section{WordCount bằng Scala}

% ------------------------------------------------------------
\subsection{Giới thiệu bài toán}

Bài toán WordCount yêu cầu đếm số lượng từ trong file \texttt{words.txt} theo chữ
cái đầu tiên, cụ thể là 8 chữ cái: \textbf{f, i, t, h, c, m, u, s} (không phân
biệt hoa thường). Chương trình cần thỏa mãn các ràng buộc sau:

\begin{itemize}
  \item Đầu vào: file \texttt{words.txt} trên HDFS tại đường dẫn
        \texttt{hdfs://localhost:9000/hcmus/23120099/words.txt}
  \item Phân tách từ bằng dấu cách (không dùng biểu thức chính quy)
  \item Chuyển tất cả từ về chữ thường trước khi kiểm tra
  \item Kết quả đầu ra: file \texttt{results.txt} có định dạng TSV (tab-separated),
        mỗi dòng gồm một chữ cái và số đếm tương ứng, tổng cộng 8 dòng
  \item Kết quả được sắp xếp theo thứ tự bảng chữ cái
\end{itemize}

Tôi sử dụng Apache Spark với ngôn ngữ Scala để thực hiện bài toán này vì Spark cung
cấp API RDD (Resilient Distributed Dataset) mạnh mẽ, phù hợp cho xử lý dữ liệu
phân tán quy mô lớn.

% ------------------------------------------------------------
\subsection{Mã nguồn Scala (WordCount.scala)}

Dưới đây là toàn bộ mã nguồn chương trình WordCount được viết bằng Scala, sử dụng
Apache Spark:

\begin{minted}[linenos]{scala}
// Nhập các thư viện cần thiết của Apache Spark
import org.apache.spark.{SparkConf, SparkContext}

// Đối tượng chính của chương trình
object WordCount {

  // Danh sách 8 chữ cái cần đếm theo yêu cầu bài lab
  val TARGET_LETTERS: Set[Char] = Set('f', 'i', 't', 'h', 'c', 'm', 'u', 's')

  def main(args: Array[String]): Unit = {

    // Bước 1: Khởi tạo cấu hình Spark với tên ứng dụng
    val conf = new SparkConf()
      .setAppName("WordCount_23120099")   // Tên hiển thị trên YARN/Spark UI
      .setMaster("local[*]")             // Chạy local, dùng toàn bộ CPU core

    // Bước 2: Tạo SparkContext — điểm khởi đầu của mọi Spark job
    val sc = new SparkContext(conf)
    sc.setLogLevel("WARN")               // Giảm log spam để dễ đọc output

    // Bước 3: Đọc file words.txt từ HDFS
    // Hadoop đã cấu hình ở localhost:9000 theo core-site.xml
    val inputPath = "hdfs://localhost:9000/hcmus/23120099/words.txt"
    val lines = sc.textFile(inputPath)   // Mỗi phần tử RDD là một dòng văn bản

    // Bước 4: Xử lý từng dòng — tách từ bằng dấu cách (không dùng regex)
    val words = lines.flatMap { line =>
      line.split(" ")                    // Phân tách bằng ký tự cách đơn giản
    }

    // Bước 5: Lọc bỏ các token rỗng phát sinh sau khi split
    val nonEmptyWords = words.filter { word =>
      word.nonEmpty                      // Loại bỏ chuỗi rỗng ""
    }

    // Bước 6: Chuyển tất cả từ về chữ thường để so sánh không phân biệt hoa thường
    val lowerWords = nonEmptyWords.map { word =>
      word.toLowerCase                   // "Hello" -> "hello", "THE" -> "the"
    }

    // Bước 7: Chỉ giữ lại các từ bắt đầu bằng một trong 8 chữ cái mục tiêu
    val filteredWords = lowerWords.filter { word =>
      word.nonEmpty && TARGET_LETTERS.contains(word.charAt(0))
      // Kiểm tra ký tự đầu tiên có thuộc tập {f, i, t, h, c, m, u, s} không
    }

    // Bước 8: Ánh xạ mỗi từ thành cặp (chữ_cái_đầu, 1) để đếm
    val letterOnePairs = filteredWords.map { word =>
      (word.charAt(0).toString, 1)       // Ví dụ: "fish" -> ("f", 1)
    }

    // Bước 9: Gộp (reduce) theo key — cộng dồn số 1 của cùng một chữ cái
    val letterCounts = letterOnePairs.reduceByKey { (count1, count2) =>
      count1 + count2                    // Ví dụ: ("f", 1), ("f", 1) -> ("f", 2)
    }

    // Bước 10: Sắp xếp kết quả theo chữ cái (key) tăng dần để output có thứ tự
    val sortedCounts = letterCounts.sortByKey(ascending = true)

    // Bước 11: Chuyển đổi sang định dạng TSV (tab-separated values)
    // Mỗi dòng: "chữ_cái\tsố_đếm" theo yêu cầu output của bài
    val tsvOutput = sortedCounts.map { case (letter, count) =>
      s"$letter\t$count"                 // Nối bằng ký tự tab '\t'
    }

    // Bước 12: Ghi kết quả ra file results.txt trên hệ thống file cục bộ
    // Dùng coalesce(1) để ghi vào một file duy nhất thay vì nhiều part file
    val outputPath = "results"
    tsvOutput.coalesce(1).saveAsTextFile(outputPath)
    // Sau khi chạy, kết quả nằm ở results/part-00000

    // Bước 13: In kết quả ra console để kiểm tra nhanh
    println("=== Kết quả WordCount ===")
    sortedCounts.collect().foreach { case (letter, count) =>
      println(s"$letter\t$count")        // Hiển thị từng dòng kết quả
    }

    // Bước 14: Dừng SparkContext để giải phóng tài nguyên
    sc.stop()
  }
}
\end{minted}

Giải thích các thành phần chính của chương trình:

\begin{itemize}
  \item \mintinline{scala}{SparkConf} và \mintinline{scala}{SparkContext}: hai đối
        tượng cốt lõi để khởi tạo môi trường Spark. \mintinline{scala}{SparkContext}
        là cổng vào duy nhất để tương tác với cluster.

  \item \mintinline{scala}{sc.textFile(inputPath)}: đọc file từ HDFS và tạo một RDD
        trong đó mỗi phần tử là một dòng văn bản. Spark tự động xử lý việc đọc phân
        tán từ các DataNode.

  \item \mintinline{scala}{flatMap}: ánh xạ mỗi dòng thành nhiều từ và làm phẳng
        (flatten) kết quả thành một RDD duy nhất chứa tất cả các từ.

  \item \mintinline{scala}{reduceByKey}: thực hiện tổng hợp phân tán — gộp tất cả
        các cặp có cùng key (chữ cái đầu) bằng hàm cộng, tương đương với bước
        \textit{Reduce} trong mô hình MapReduce.

  \item \mintinline{scala}{sortByKey}: sắp xếp RDD theo key theo thứ tự từ điển,
        đảm bảo output có thứ tự nhất quán.

  \item \mintinline{scala}{coalesce(1)}: hợp nhất tất cả partition thành một file
        duy nhất trước khi ghi, tránh tình trạng kết quả bị phân mảnh thành nhiều
        \texttt{part-XXXXX} file.
\end{itemize}

% ------------------------------------------------------------
\subsection{Hướng dẫn biên dịch và chạy}

Dự án được tổ chức theo cấu trúc SBT (Scala Build Tool) tiêu chuẩn. Để biên dịch
và đóng gói chương trình, tôi chạy lệnh sau trong thư mục gốc của dự án:

\begin{minted}{bash}
sbt package
\end{minted}

Lệnh này biên dịch toàn bộ mã nguồn Scala và tạo ra file JAR tại
\texttt{target/scala-2.12/wordcount\_2.12-1.0.jar}. Sau khi có file JAR, tôi submit
job lên Spark:

\begin{minted}{bash}
spark-submit \
  --class WordCount \
  --master local[*] \
  target/scala-2.12/wordcount_2.12-1.0.jar
\end{minted}

Spark sẽ đọc file \texttt{words.txt} trực tiếp từ HDFS thông qua đường dẫn
\texttt{hdfs://localhost:9000/hcmus/23120099/words.txt} đã được hardcode trong mã
nguồn. Quá trình này sử dụng thư viện Hadoop client tích hợp sẵn trong Spark để
giao tiếp với NameNode và DataNode. Với chế độ \texttt{local[*]}, Spark sử dụng
tất cả CPU core của máy để thực thi job mà không cần cluster YARN riêng biệt.

% Ảnh chụp màn hình kết quả spark-submit
\begin{figure}[H]
  \centering
  \includegraphics[width=0.9\textwidth]{spark_submit.png}
  \caption{Chạy spark-submit WordCount thành công}
  \label{fig:spark-submit}
\end{figure}

% ------------------------------------------------------------
\subsection{Kết quả}

Sau khi chạy xong, kết quả được lưu vào file \texttt{results/part-00000} với định
dạng tab-separated. Nội dung file \texttt{results.txt} thu được:

\begin{minted}{text}
c	38934
f	15870
h	18662
i	15220
m	25194
s	50571
t	25223
u	23790
\end{minted}

% Ảnh chụp màn hình cat results.txt
\begin{figure}[H]
  \centering
  \includegraphics[width=0.9\textwidth]{result.png}
  \caption{Nội dung file results.txt sau khi chạy WordCount}
  \label{fig:results}
\end{figure}

% ------------------------------------------------------------
\subsection{Kiểm thử}

Để xác minh tính đúng đắn của chương trình trước khi chạy trên toàn bộ
\texttt{words.txt}, tôi tự kiểm thử với một tập dữ liệu nhỏ gồm 10 từ bao phủ
đủ 8 chữ cái mục tiêu:

\textbf{Dữ liệu kiểm thử (\texttt{test\_input.txt}):}
\begin{minted}{text}
fish Imagine the House cat Mango umbrella Sun Fire hello
\end{minted}

\textbf{Kết quả tính tay (không phân biệt hoa thường):}
\begin{itemize}
  \item \textbf{c}: cat $\rightarrow$ 1
  \item \textbf{f}: fish, Fire $\rightarrow$ 2
  \item \textbf{h}: House, hello $\rightarrow$ 2
  \item \textbf{i}: Imagine $\rightarrow$ 1
  \item \textbf{m}: Mango $\rightarrow$ 1
  \item \textbf{s}: Sun $\rightarrow$ 1
  \item \textbf{t}: the $\rightarrow$ 1
  \item \textbf{u}: umbrella $\rightarrow$ 1
\end{itemize}

\textbf{Kết quả mong đợi (TSV):}
\begin{minted}{text}
c	1
f	2
h	2
i	1
m	1
s	1
t	1
u	1
\end{minted}

\textbf{Kết quả thực tế từ chương trình:} Kết quả từ chạy WordCount trên file kiểm
thử khớp hoàn toàn với kết quả tính tay, xác nhận rằng chương trình hoạt động đúng
theo logic MapReduce: đọc từ HDFS, phân tách theo dấu cách, chuyển chữ thường, lọc
8 chữ cái mục tiêu, đếm, sắp xếp, và xuất định dạng TSV.
